% GENERATED FILE. Edit canonical lessons, then run npm run generate:books.
% canonical-source-hash: fnv1a64:e1703a8e47485f84
% canonical-lessons: AR-C365-haqiqa, AR-C365-waqi, AR-C365-khayal, AR-C365-fannada, AR-C365-iqtanaa, AR-R365-review-358-365

\chapter{Fact and fancy}
\label{ch:ar-fact-and-fancy}
\hlchaptermodality{365}

\begin{chapteropening}
\textbf{By the end of this chapter:} \emph{I can talk about fact, reality and imagination, and say that someone refuted a claim or was convinced.}

Complete the last lesson of chapter 365: I can talk about fact, reality and imagination, and say that someone refuted a claim or was convinced.
\end{chapteropening}

\section[\texorpdfstring{\ar{حقيقة} (ḥaqīqa)}{حقيقة (ḥaqīqa)}]{\ar{حقيقة} (ḥaqīqa) — a fact}

\label{lesson:AR-C365-haqiqa}



\begin{quote}
Before the new one: say the Arabic for to specify, then the Arabic for to distinguish.
\end{quote}

\subsection*{You'll want to know: \ar{حقيقة}}
\textbf{\ar{حقيقة}} — \emph{ḥaqīqa} — \textquotedblleft{}a fact\textquotedblright{}.

\begin{quote}
\ar{هذه} \ar{حقيقة} \ar{لا} \ar{يعرفها} \ar{كثير} \ar{من} \ar{الناس}.
\end{quote}

\begin{quote}
\emph{hādhihi ḥaqīqa lā yaʿrifuhā kathīr min an-nās.}
\end{quote}

\begin{quote}
This is a fact that many people do not know.
\end{quote}

\ar{في} \ar{الحقيقة}, \emph{fī al-ḥaqīqa}, in fact.

\begin{cousinweb}
\textbf{\ar{حقيقة}} (\emph{ḥaqīqa}) is built on the native Arabic root \emph{ḥ-q-q}, \textquotedblleft{}to be true, to be due\textquotedblright{} (Lane). A \emph{ḥaqīqa} is what is truly so: a fact, the truth.
\end{cousinweb}

\begin{grammarlens}[title={What you've built}]
One more word for giving reasons and opinions.
\end{grammarlens}

\subsection*{Your turn}
\begin{itemize}
\raggedright
  \item \emph{Say it:} \emph{ḥaqīqa}
  \item \emph{Say it:} \emph{ḥaqīqa}, once more
  \item \emph{Say it:} \emph{hādhihi ḥaqīqa lā yaʿrifuhā kathīr min an-nās.} — This is a fact that many people do not know.
  \item \emph{Recall:} say the Arabic for to specify, then the Arabic for to distinguish, then say \emph{ḥaqīqa} again
\end{itemize}

\begin{tcolorbox}[breakable,colback=teal!4,colframe=teal!35!black,title={Before you move on}]
What does \ar{حقيقة} mean? (\textquotedblleft{}A fact\textquotedblright{}.) Say it once more.
\end{tcolorbox}
\section[\texorpdfstring{\ar{واقع} (wāqiʿ)}{واقع (wāqiʿ)}]{\ar{واقع} (wāqiʿ) — reality}

\label{lesson:AR-C365-waqi}



\begin{quote}
Before the new one: say the Arabic for to distinguish, then the Arabic for a fact.
\end{quote}

\subsection*{You'll want to know: \ar{واقع}}
\textbf{\ar{واقع}} — \emph{wāqiʿ} — \textquotedblleft{}reality\textquotedblright{}.

\begin{quote}
\ar{الواقع} \ar{مختلف} \ar{عن} \ar{الأفلام}.
\end{quote}

\begin{quote}
\emph{al-wāqiʿ mukhtalif ʿan al-aflām.}
\end{quote}

\begin{quote}
Reality is different from the films.
\end{quote}

\ar{في} \ar{الواقع}, \emph{fī al-wāqiʿ}, in reality, actually.

\begin{cousinweb}
\textbf{\ar{واقع}} (\emph{wāqiʿ}) is built on the native Arabic root \emph{w-q-ʿ}, \textquotedblleft{}to fall, to happen\textquotedblright{} (Lane). \emph{Wāqiʿ} is the active participle, what has fallen, what happens: reality.
\end{cousinweb}

\begin{grammarlens}[title={What you've built}]
One more word for giving reasons and opinions.
\end{grammarlens}

\subsection*{Your turn}
\begin{itemize}
\raggedright
  \item \emph{Say it:} \emph{wāqiʿ}
  \item \emph{Say it:} \emph{wāqiʿ}, once more
  \item \emph{Say it:} \emph{al-wāqiʿ mukhtalif ʿan al-aflām.} — Reality is different from the films.
  \item \emph{Recall:} say the Arabic for to distinguish, then the Arabic for a fact, then say \emph{wāqiʿ} again
\end{itemize}

\begin{tcolorbox}[breakable,colback=teal!4,colframe=teal!35!black,title={Before you move on}]
What does \ar{واقع} mean? (\textquotedblleft{}Reality\textquotedblright{}.) Say it once more.
\end{tcolorbox}
\section[\texorpdfstring{\ar{خيال} (khayāl)}{خيال (khayāl)}]{\ar{خيال} (khayāl) — imagination}

\label{lesson:AR-C365-khayal}



\begin{quote}
Before the new one: say the Arabic for a fact, then the Arabic for reality.
\end{quote}

\subsection*{You'll want to know: \ar{خيال}}
\textbf{\ar{خيال}} — \emph{khayāl} — \textquotedblleft{}imagination\textquotedblright{}.

\begin{quote}
\ar{عند} \ar{ابنتي} \ar{خيال} \ar{واسع}.
\end{quote}

\begin{quote}
\emph{ʿinda ibnatī khayāl wāsiʿ.}
\end{quote}

\begin{quote}
My daughter has a vivid imagination.
\end{quote}

\ar{خيال} \ar{علمي}, \emph{khayāl ʿilmī}, science fiction.

\begin{cousinweb}
\textbf{\ar{خيال}} (\emph{khayāl}) is built on the native Arabic root \emph{kh-y-l}, \textquotedblleft{}to think, to fancy\textquotedblright{} (Lane). \emph{Khayāl} is an image or phantom in the mind: imagination.
\end{cousinweb}

\begin{grammarlens}[title={What you've built}]
One more word for giving reasons and opinions.
\end{grammarlens}

\subsection*{Your turn}
\begin{itemize}
\raggedright
  \item \emph{Say it:} \emph{khayāl}
  \item \emph{Say it:} \emph{khayāl}, once more
  \item \emph{Say it:} \emph{ʿinda ibnatī khayāl wāsiʿ.} — My daughter has a vivid imagination.
  \item \emph{Recall:} say the Arabic for a fact, then the Arabic for reality, then say \emph{khayāl} again
\end{itemize}

\begin{tcolorbox}[breakable,colback=teal!4,colframe=teal!35!black,title={Before you move on}]
What does \ar{خيال} mean? (\textquotedblleft{}Imagination\textquotedblright{}.) Say it once more.
\end{tcolorbox}
\section[\texorpdfstring{\ar{فند} (fannada)}{فند (fannada)}]{\ar{فند} (fannada) — to refute}

\label{lesson:AR-C365-fannada}



\begin{quote}
Before the new one: say the Arabic for reality, then the Arabic for imagination.
\end{quote}

\subsection*{You'll want to know: \ar{فند}}
\textbf{\ar{فند}} — \emph{fannada} — \textquotedblleft{}to refute\textquotedblright{}.

\begin{quote}
\ar{فند} \ar{المحامي} \ar{كلام} \ar{الشاهد}.
\end{quote}

\begin{quote}
\emph{fannada al-muḥāmī kalām ash-shāhid.}
\end{quote}

\begin{quote}
The lawyer refuted the witness's statement.
\end{quote}

Present: \ar{يفند}, \emph{yufannid}.

\begin{cousinweb}
\textbf{\ar{فند}} (\emph{fannada}) is built on the native Arabic root \emph{f-n-d}, \textquotedblleft{}to lie, to say what is untrue\textquotedblright{} (Lane). It is Form II, to show to be untrue: to refute.
\end{cousinweb}

\begin{grammarlens}[title={What you've built}]
One more word for giving reasons and opinions.
\end{grammarlens}

\subsection*{Your turn}
\begin{itemize}
\raggedright
  \item \emph{Say it:} \emph{fannada}
  \item \emph{Say it:} \emph{fannada}, once more
  \item \emph{Say it:} \emph{fannada al-muḥāmī kalām ash-shāhid.} — The lawyer refuted the witness's statement.
  \item \emph{Recall:} say the Arabic for reality, then the Arabic for imagination, then say \emph{fannada} again
\end{itemize}

\begin{tcolorbox}[breakable,colback=teal!4,colframe=teal!35!black,title={Before you move on}]
What does \ar{فند} mean? (\textquotedblleft{}To refute\textquotedblright{}.) Say it once more.
\end{tcolorbox}
\section[\texorpdfstring{\ar{اقتنع} (iqtanaʿa)}{اقتنع (iqtanaʿa)}]{\ar{اقتنع} (iqtanaʿa) — to be convinced}

\label{lesson:AR-C365-iqtanaa}



\begin{quote}
Before the new one: say the Arabic for imagination, then the Arabic for to refute.
\end{quote}

\subsection*{You'll want to know: \ar{اقتنع}}
\textbf{\ar{اقتنع}} — \emph{iqtanaʿa} — \textquotedblleft{}to be convinced\textquotedblright{}.

\begin{quote}
\ar{اقتنع} \ar{أبي} \ar{بالفكرة} \ar{بعد} \ar{ساعة}.
\end{quote}

\begin{quote}
\emph{iqtanaʿa abī bi-l-fikra baʿda sāʿa.}
\end{quote}

\begin{quote}
My father was convinced by the idea after an hour.
\end{quote}

Present: \ar{يقتنع}, \emph{yaqtaniʿ}.

It takes \ar{بـ}: \ar{اقتنع} \ar{بـ} …. Compare \ar{أقنع}, to convince.

\begin{cousinweb}
\textbf{\ar{اقتنع}} (\emph{iqtanaʿa}) is built on the native Arabic root \emph{q-n-ʿ}, \textquotedblleft{}to be content, satisfied\textquotedblright{} (Lane). It is Form VIII, to become content with: to be convinced.
\end{cousinweb}

\begin{grammarlens}[title={What you've built}]
One more word for giving reasons and opinions.
\end{grammarlens}

\subsection*{Your turn}
\begin{itemize}
\raggedright
  \item \emph{Say it:} \emph{iqtanaʿa}
  \item \emph{Say it:} \emph{iqtanaʿa}, once more
  \item \emph{Say it:} \emph{iqtanaʿa abī bi-l-fikra baʿda sāʿa.} — My father was convinced by the idea after an hour.
  \item \emph{Recall:} say the Arabic for imagination, then the Arabic for to refute, then say \emph{iqtanaʿa} again
\end{itemize}

\begin{tcolorbox}[breakable,colback=teal!4,colframe=teal!35!black,title={Before you move on}]
What does \ar{اقتنع} mean? (\textquotedblleft{}To be convinced\textquotedblright{}.) Say it once more.
\end{tcolorbox}
\section[(dialogue)]{Reasons and arguments — a review}

\label{lesson:AR-R365-review-358-365}



\begin{quote}
Say the Arabic for an argument and a principle.
\end{quote}

\subsection*{Your turn}
Say these aloud:

\begin{itemize}
\raggedright
  \item an argument, a principle, a proof, to prove, to convince
  \item an impression, a perspective, to assume, to suggest, convincing
  \item to approve, to object, consensus, to argue, a controversy
  \item a public issue, a matter, a context, to stir up, to wonder
  \item logic, a fallacy, the mind, to deduce, a contradiction
  \item a theory, a hypothesis, a concept, a phenomenon, to interpret
  \item to analyse, an element, a criterion, to specify, to distinguish
  \item a fact, reality, imagination, to refute, to be convinced
\end{itemize}

\begin{tcolorbox}[breakable,colback=teal!4,colframe=teal!35!black,title={Before you move on}]
An argument? (\textbf{\ar{حجة}}, \emph{ḥujja}.) A principle? (\textbf{\ar{مبدأ}}, \emph{mabdaʾ}.) A proof? (\textbf{\ar{برهان}}, \emph{burhān}.) To prove? (\textbf{\ar{أثبت}}, \emph{athbata}.) To convince? (\textbf{\ar{أقنع}}, \emph{aqnaʿa}.) An impression? (\textbf{\ar{انطباع}}, \emph{inṭibāʿ}.) A perspective? (\textbf{\ar{منظور}}, \emph{manẓūr}.) To assume? (\textbf{\ar{افترض}}, \emph{iftaraḍa}.) To suggest? (\textbf{\ar{اقترح}}, \emph{iqtaraḥa}.) Convincing? (\textbf{\ar{مقنع}}, \emph{muqniʿ}.) To approve? (\textbf{\ar{وافق}}, \emph{wāfaqa}.) To object? (\textbf{\ar{اعترض}}, \emph{iʿtaraḍa}.) Consensus? (\textbf{\ar{إجماع}}, \emph{ijmāʿ}.) To argue? (\textbf{\ar{جادل}}, \emph{jādala}.) A controversy? (\textbf{\ar{جدل}}, \emph{jadal}.) A public issue? (\textbf{\ar{قضية}}, \emph{qaḍiyya}.) A matter? (\textbf{\ar{مسألة}}, \emph{masʾala}.) A context? (\textbf{\ar{سياق}}, \emph{siyāq}.) To stir up? (\textbf{\ar{أثار}}, \emph{athāra}.) To wonder? (\textbf{\ar{تساءل}}, \emph{tasāʾala}.) Logic? (\textbf{\ar{منطق}}, \emph{manṭiq}.) A fallacy? (\textbf{\ar{مغالطة}}, \emph{mughālaṭa}.) The mind? (\textbf{\ar{عقل}}, \emph{ʿaql}.) To deduce? (\textbf{\ar{استنتج}}, \emph{istantaja}.) A contradiction? (\textbf{\ar{تناقض}}, \emph{tanāquḍ}.) A theory? (\textbf{\ar{نظرية}}, \emph{naẓariyya}.) A hypothesis? (\textbf{\ar{فرضية}}, \emph{farḍiyya}.) A concept? (\textbf{\ar{مفهوم}}, \emph{mafhūm}.) A phenomenon? (\textbf{\ar{ظاهرة}}, \emph{ẓāhira}.) To interpret? (\textbf{\ar{فسر}}, \emph{fassara}.) To analyse? (\textbf{\ar{حلل}}, \emph{ḥallala}.) An element? (\textbf{\ar{عنصر}}, \emph{ʿunṣur}.) A criterion? (\textbf{\ar{معيار}}, \emph{miʿyār}.) To specify? (\textbf{\ar{حدد}}, \emph{ḥaddada}.) To distinguish? (\textbf{\ar{ميز}}, \emph{mayyaza}.) A fact? (\textbf{\ar{حقيقة}}, \emph{ḥaqīqa}.) Reality? (\textbf{\ar{واقع}}, \emph{wāqiʿ}.) Imagination? (\textbf{\ar{خيال}}, \emph{khayāl}.) To refute? (\textbf{\ar{فند}}, \emph{fannada}.) To be convinced? (\textbf{\ar{اقتنع}}, \emph{iqtanaʿa}.)
\end{tcolorbox}
